% 解説プリントのコメント（saiten が書く。./tools/kaisetsu.ps1 がこのファイルからプリントを作る）
%% matome
\begin{daiji}
  \item エルミート性 $\int\phi^{*}(\hat{A}\psi)\dd x=\int(\hat{A}\phi)^{*}\psi\dd x$ は積分どうしの等式で、被積分関数どうしは等しくない。固有値の実数性は積分を付けたまま示し、最後に $\int\abs{\psi}^{2}\dd x>0$ で割る（問題7 (2)）。
  \item 「どの段階で成り立たないか」「基底によらないのはなぜか」と聞かれたら、式で確かめた後に理由を1文書く。非エルミート行列では $B^{\dagger}\vect{v}$ が $\vect{v}$ の定数倍にならない（問題7 (4)）。縮退した固有ベクトルの1次結合も同じ固有値の固有ベクトルなので直交化してよい（問題8 (2)）。縮退した部分の係数の2乗和は固有空間への射影の長さの2乗である（問題8 (4)）。
\end{daiji}
\begin{maemo}
  \item 根拠の1文の書き落とし：量子力学1-2 問題1 (1)（2026-10-14、位相一定の根拠なし）、量子力学1-4 問題3 (2)（2026-10-28、$x=-1$ が区間外である根拠なし）でも、結論に必要な根拠を書いていなかった。今回は問題7 (2)(4)、問題8 (2)(4) で同じことが起きた。答えを書いたら「なぜそう言えるか」を1行足してから次の問いへ進む。
  \item 複素共役の書き方：量子力学1-3 問題2（2026-10-18・2026-10-21、複素定数の $\abs{A}^{2}$ を $A^{2}$）、問題5（2026-10-18、$R$ を $R^{*}$ にしなかった）。今回も問題7 (1)〜(3) で $(\hat{A}\psi)^{*}$ を $\hat{A}^{\dagger}\psi^{*}$ と書いていた（最後の式は正しいので減点していない）。共役は括弧ごと $(\hat{A}\psi)^{*}$ と書き、$\hat{A}^{\dagger}$（エルミート共役）と $^{*}$（複素共役）を混ぜない。
\end{maemo}
%% mondai 7
\begin{kotaeran}
  \machigai{2}{○}{記述} $\hat{A}\psi=\lambda\psi$ から $\psi^{*}\hat{A}\psi=\lambda\psi^{*}\psi=(\hat{A}^{\dagger}\psi^{*})\psi=\lambda^{*}\psi^{*}\psi$、$\therefore\lambda=\lambda^{*}$ と書いている。方針（2通りに計算して比べる）は正しいが、積分記号がすべて落ちている。$\psi^{*}(\hat{A}\psi)=(\hat{A}\psi)^{*}\psi$ は各点では成り立たず（例：$\hat{p}$ では左辺と右辺で符号が違う）、エルミート性が保証するのは積分が等しいことだけである。また、$\lambda\psi^{*}\psi=\lambda^{*}\psi^{*}\psi$ から $\lambda=\lambda^{*}$ とするには、$\psi\neq0$ なので $\int\abs{\psi}^{2}\dd x>0$ であり、それで割ってよいことが要る。正しくは、(1) の $\psi$ に固有関数を使って $\int\psi^{*}\hat{A}\psi\dd x=\lambda\int\abs{\psi}^{2}\dd x$ と書き、左辺が (1) より実数、$\int\abs{\psi}^{2}\dd x$ が正の実数なので $\lambda$ は実数、とする。

  \machigai{4}{○}{記述} 固有値 $1$、$2$、固有ベクトル $(1,0)^{\mathrm{T}}$、$(1,1)^{\mathrm{T}}/\sqrt{2}$、内積 $1/\sqrt{2}\neq0$ は正しい。(3) の書き直しとして $\vect{v}_{2}^{\dagger}B\vect{v}_{1}=\vect{v}_{2}^{\dagger}\vect{v}_{1}$ と $(B^{\dagger}\vect{v}_{2})^{\dagger}\vect{v}_{1}=2\vect{v}_{2}^{\dagger}\vect{v}_{1}$ を並べ、後者に波線を引いているが、「この段階が成り立たない」ことも、その理由も書いていない。「述べよ」の問いなので、ここで $B^{\dagger}\vect{v}_{2}=2\vect{v}_{2}$（つまり $B^{\dagger}=B$）を使っているが $B$ はエルミートでなく、実際 $B^{\dagger}\vect{v}_{2}=(1,3)^{\mathrm{T}}/\sqrt{2}$ は $\vect{v}_{2}$ の定数倍でない、と書く必要がある。

  問題とは別に (1)(3) の書き方について：$\int(\hat{A}^{\dagger}\psi^{*})\psi\dd x=\int\psi\hat{A}\psi^{*}\dd x$ のように、$(\hat{A}\psi)^{*}$ を $\hat{A}^{\dagger}\psi^{*}$ や $\hat{A}\psi^{*}$ と書いている。文字どおりには「$\psi^{*}$ に $\hat{A}$ を作用させたもの」で、$\hat{p}$ なら $\hat{p}(\psi^{*})=-(\hat{p}\psi)^{*}$ と符号まで変わる。最後の式は正しく、意図は読み取れるので減点していないが、括弧ごと $(\hat{A}\psi)^{*}$ と書く。(3) の途中の $\int\phi_{m}^{*}u_{n}\dd x$ は $\phi_{n}$ の書き間違い（次の行で直っている）。
\end{kotaeran}
\begin{rikai}
  \item エルミート演算子の定義 $\int\phi^{*}(\hat{A}\psi)\dd x=\int(\hat{A}\phi)^{*}\psi\dd x$ は、部分積分で演算子を左の関数へ移したときの等式である。だから積分記号を外した形では使えない。$\hat{p}$ の場合、部分積分の境界の項が 0 になって初めて成り立つ。
  \item 固有値の実数性の証明の要は2つ：(1) の「$\int\psi^{*}\hat{A}\psi\dd x$ は実数」と、$\int\abs{\psi}^{2}\dd x$ が 0 でない正の実数であること。後者が無いと $\lambda$ について割れない。
  \item (3) の証明でエルミート性を使うのは「$\hat{A}$ を左に移して、左の固有関数の固有値で置き換える」ところだけである。行列では $\vect{u}^{\dagger}B\vect{v}=(B^{\dagger}\vect{u})^{\dagger}\vect{v}$ の $B^{\dagger}\vect{u}$ が $\vect{u}$ の定数倍かどうかを、実際に掛け算して確かめるとよい。
\end{rikai}
%% mondai 8
\begin{kotaeran}
  \machigai{2}{○}{記述} $\vect{u}_{a}$、$\vect{u}_{b}$ が固有値 $1$ の固有ベクトルであることの確認と、Gram--Schmidt の結果 $\vect{e}_{1}=(1,-1,0)^{\mathrm{T}}/\sqrt{2}$、$\vect{e}_{2}=(1,1,-2)^{\mathrm{T}}/\sqrt{6}$ は正しい。抜けているのは2点。(a) 問いの「互いに直交しないことを確かめよ」に当たる $\vect{u}_{a}^{\dagger}\vect{u}_{b}=1\neq0$ が書かれていない（射影の係数 $1/2$ から読み取れるだけ）。(b) 直交化してできた $\vect{e}_{2}$ が固有ベクトルであることの根拠、つまり「同じ固有値の固有ベクトルの1次結合は、また同じ固有値の固有ベクトル」が無い。また分母を $\abs{\vect{u}_{2}-(\vect{u}_{2}\cdot\vect{u}_{1})\vect{u}_{1}}$ と書いているが、正しくは $\vect{e}_{1}$ 方向の成分を引いた $\abs{\vect{u}_{2}-(\vect{u}_{2}\cdot\vect{e}_{1})\vect{e}_{1}}$ である（値 $\sqrt{2/3}$ は正しい）。

  \machigai{4}{○}{記述} 連立方程式から $c_{1}'=1/\sqrt{2}$、$c_{2}'=1/\sqrt{6}$ を出し、$c_{1}'^{2}+c_{2}'^{2}=2/3=c_{1}^{2}+c_{2}^{2}$ を確かめたところまでは正しい。足りないのは、なぜ基底を取り替えても変わらないかの1文である。縮退した部分の係数の2乗和は、$\vect{v}$ を固有値 $1$ の固有空間へ射影したベクトルの長さの2乗で、$\abs{\vect{v}}^{2}-\abs{c_{0}}^{2}=1-1/3=2/3$ と基底によらずに決まる。
\end{kotaeran}
\begin{rikai}
  \item 縮退した固有値 $\lambda$ では、$A\vect{u}_{a}=\lambda\vect{u}_{a}$、$A\vect{u}_{b}=\lambda\vect{u}_{b}$ から $A(\alpha\vect{u}_{a}+\beta\vect{u}_{b})=\lambda(\alpha\vect{u}_{a}+\beta\vect{u}_{b})$ となる。だから固有空間の中で Gram--Schmidt を行っても固有ベクトルのままである。これが直交化してよい理由になる。
  \item 規格直交基底での係数の2乗和は、そのベクトルの長さの2乗（射影の長さの2乗）である。固有空間が決まれば射影は決まるので、その中の基底の選び方によらない。
  \item 展開係数は連立方程式でも求まるが、規格直交基底なら $c_{j}=\vect{e}_{j}^{\dagger}\vect{v}$ と内積1回で求まる。連立方程式で解いたときは、内積で検算するとよい。
\end{rikai}
